\documentclass[11pt]{article}

\usepackage[margin=1.1in]{geometry}
\usepackage{amsmath,amssymb}
\usepackage{hyperref}

\title{BEM Solver}
\author{}
\date{}

\begin{document}
\maketitle

\begin{abstract}
These notes describe a cylindrically symmetric Boundary Element Method (BEM) that is used to solve for the electrostatic environment around a nano-structured photocathode. We solve for both the static accelerating fields and also the ``image charge'' fields induced on any emitted electrons. Though the problem is generally three-dimensional, we show that expanding in angular harmonics can reduce it to independent two-dimensional problems, one for each
azimuthal harmonic $m$.  This makes it possible to treat off-axis point charges
without ever introducing an angular mesh, and makes an overall rotation of the point sources an exact symmetry of the method. Finally, we explain how for certain nano-geometries, it is possible to use inversion symmetry to handle the boundary condition on the infinite grounded plane, without needing a mesh on that boundary.
\end{abstract}

\section{What problem are we solving?}

We model the nano-structured photocathode as a perfect conductor with surface $S$, surrounded by a vacuum region denoted by $\Omega$. We only treat the case of a single conductor, so we choose the potential on the conductor's surface to be zero.
An externally applied potential
$\Phi_{\mathrm{ext}}$ is added inside $\Omega$ (e.g. a constant accelerating field, or the Coulomb field from an emitted electron), and the
conductor responds by producing an induced potential $\Phi$ such that
\begin{equation}
  \nabla^2 \Phi = 0 \ \text{ in } \Omega,
  \qquad
  \Phi = -\Phi_{\mathrm{ext}} \ \text{ on } S,
  \qquad
  \Phi \to 0 \ \text{ at infinity}.
  \label{eq:dirichlet}
\end{equation}
The total potential $\Phi_{\mathrm{ext}}+\Phi$ therefore vanishes on the conductor and goes asymptotically to any applied field far from the cathode.

Boundary Element Methods do not solve Laplace's equation on a three-dimensional volume.
Instead, we write the field in terms of an unknown surface charge density $\sigma$ on
the conductor,
\begin{equation}
  \Phi(\mathbf{x}) = \frac{1}{4\pi\varepsilon_0}
  \int_S \frac{\sigma(\mathbf{y})}{|\mathbf{x}-\mathbf{y}|}
  \, dS(\mathbf{y}).
  \label{eq:single-layer}
\end{equation}
We then choose $\sigma$ so that this potential satisfies the boundary condition
in Eq.~\eqref{eq:dirichlet}.  Numerically, the surface is broken up into a finite number of segments that allow for an approximate solution to the problem.

The same machinery is used for two physically different excitations.

\begin{description}
\item[Applied field.] For a uniform field along $z$,
  $\Phi_{\mathrm{ext}}=-E_z z$.  The induced surface charge tells us how the
  conductor distorts and screens that field around the geometric feature.

\item[Point charges / image response.] Here $\Phi_{\mathrm{ext}}$ is the
  Coulomb potential of one or more charges in the vacuum.  Solving for the
  induced surface charge gives the induced image field.
\end{description}

The distinction matters because the applied field case is cylindrically symmetric, while an
off-axis point charge is not.  The cathode geometry itself may still be cylindrically symmetric in both cases, but the excitation breaks that symmetry.

\section{Using rotational symmetry without assuming a cylindrically symmetric source}

Suppose the conductor is a surface of revolution about $\hat z$, defined by a curve $\Gamma$ in $(\rho,z)$. Write the surface charge density as a sum over angular Fourier components,
\begin{equation}
  \sigma(s,\phi) = \sum_{m \ge 0}
  \Big[ \sigma^c_m(s)\cos m\phi + \sigma^s_m(s)\sin m\phi \Big],
  \qquad s \ \text{along } \Gamma.
\end{equation}
and let the boundary conditions of the potential be specified by similarly defined Fourier coefficients $g^{c,s}_m$. Orthogonality of the Fourier components can be used to show that,
\begin{equation}
  \int_\Gamma K_m(\rho,z;\rho',z')\,
  \sigma^{c,s}_m(s')\, \rho'\, ds'
  = g^{c,s}_m(s),
  \qquad m = 0,1,2,\dots,
  \label{eq:per-mode}
\end{equation}
where
\begin{equation}
  K_m(\rho,z;\rho',z')
  = \frac{Q_{m-\frac12}(\chi)}{2\pi \varepsilon_0 \sqrt{\rho\rho'}},
  \label{eq:ring-kernel}
\end{equation}
$Q_{m-\frac12}$ is a half-integer Legendre function of the second kind, and $\chi$ is defined as
\begin{equation}
  \chi = \frac{\rho^2 + \rho'^2 + (z-z')^2}{2\rho\rho'} \ > 1.
\end{equation}
That is, each angular mode of the source potential only excites that particular angular mode in the induced potential, and the problem for each mode is purely in the $(\rho,z)$ plane.

Two practical facts are important here.  First, the modes are completely
uncoupled.  With $P$ boundary nodes, we solve one $P\times P$ system $A_m$ for
each $m$, not one giant $(MP)\times(MP)$ system.  Second, the same matrix $A_m$
applies to both the cosine and sine coefficients.  The geometry and kernel do
not distinguish between them; only the excitation does.

Conveniently, a point source can be expanded in the same basis using the same functions,
\begin{equation}
  \frac{1}{|\mathbf{x}-\mathbf{x}_0|}
  = \frac{1}{\pi\sqrt{\rho\rho_0}}
    \sum_{m\ge 0} (2-\delta_{m0})\, Q_{m-\frac12}(\chi_0)\,
    \cos\!\big(m(\phi-\phi_0)\big),
  \label{eq:point-expansion}
\end{equation}
where $\chi_0$ is formed from the field point $(\rho,z)$ and source position
$(\rho_0,z_0)$.  These Fourier coefficients supply the right-hand sides $g_m^{c,s}$ for the
image-response solve.

A useful way to think about this is that the azimuthal coordinate is treated
analytically rather than meshed.  There is no $\phi$ grid anywhere in the
solver.  The source azimuth enters exactly through the Fourier factors, and the
only azimuthal approximation is that in practice the series is stopped at some
finite $m=M$.

\subsection{Evaluating the kernels near coincidence}

Write $\eta=\operatorname{arccosh}\chi$. The unwanted solution in an upward
Legendre recurrence grows relative to $Q$ roughly as $e^{2m\eta}$. Thus
upward recurrence is unsuitable at large $m\eta$, but can be used near
coincidence for a finite range of modes. The implementation uses the
elliptic-integral seeds $Q_{-1/2}$ and $Q_{1/2}$ and recurs upward when
$\max(M,1)\eta\leq1$. Elsewhere it uses downward ratio recurrence.
The switch applies separately to each field/source pair, including on the
GPU. Values and derivatives on both sides of the switch are checked
against high-precision Legendre functions.

The small complementary elliptic parameter is computed directly as
$p=(\chi-1)/(\chi+1)$, rather than by subtracting an already-rounded
parameter from one. Likewise, derivative denominators use
$(\chi-1)(\chi+1)$ instead of $\chi^2-1$.

For the static field, nearby meridional segments require more quadrature
resolution than distant ones. With $s_c$ the closest point on a segment
and $d$ its distance from the query, the substitution
$s=s_c+d\sinh v$ resolves the narrow peak with a number of quadrature
points growing only logarithmically with segment length divided by $d$.
This integrates the actual surface density; replacing a constant density
over the whole surface by a local planar field is not an exact identity.
Queries exactly on the piecewise-linear surface use a small exterior
displacement proportional to the local element length. At polygon
corners a unique finite field need not exist, so profile refinement
remains necessary when interpreting surface fields.

\section{What does truncating the azimuthal series actually do?}

Stopping the point-source expansion at $m\le M$ is the main approximation in
the azimuthal direction.  It removes fine angular structure, but it does
\emph{not} introduce preferred absolute angles.

To see why, rotate every source by the same angle $\alpha$.  For each fixed
mode $m$, the cosine and sine components transform as
\begin{equation}
  \begin{pmatrix} g^c_m \\ g^s_m \end{pmatrix}
  \longmapsto
  \begin{pmatrix} \cos m\alpha & -\sin m\alpha \\
                  \sin m\alpha & \cos m\alpha \end{pmatrix}
  \begin{pmatrix} g^c_m \\ g^s_m \end{pmatrix}.
  \label{eq:mode-rotation}
\end{equation}
So a physical rotation by $\alpha$ simply rotates the cosine/sine pair of mode
$m$ by the angle $m\alpha$.  Crucially, it never mixes one value of $m$ with
another.

The same matrix $A_m$ acts on both members of the cosine/sine pair, and $A_m$
does not depend on the absolute source azimuth.  Therefore the solved surface
charge coefficients rotate in exactly the same way as the boundary data.  If
the field point is rotated along with the sources, the reconstructed field is
just the rigidly rotated version of the original field.

This remains true after truncation because truncation removes complete modes.
Keeping $m=0,\ldots,M$ still leaves each retained mode with its full cosine/sine
pair and therefore with its exact rotational transformation law.  In that
sense, rotational covariance is exact at every $M$: changing the absolute
orientation of the same physical configuration does not change the numerical
answer except for rotating the vector components. For several point charges, this means the computed forces depend on azimuths only through relative angles $\phi_i-\phi_j$.

What truncation \emph{does} cost is angular resolution.  A charge very close to
the conductor produces a narrow peak in the induced surface charge, and a
narrower peak requires larger $m$ to represent.  The important angular scale is
therefore set mainly by the charge's standoff from the surface, not by the
azimuthal separation between two charges.

There is one coordinate special case at $\rho\to0$.  The azimuth becomes
undefined on the axis and $\chi\to\infty$ in Eq.~\eqref{eq:ring-kernel}.  This
is a cylindrical-coordinate singularity rather than a failure of the mode
expansion, and numerically it is handled with a floor on $\rho$.  A source
exactly on the axis is even simpler: the physical configuration itself becomes
axisymmetric, only $m=0$ is excited, and the transverse field at the source
vanishes by symmetry.  This should not be confused with evaluating the field
\emph{on} the axis from an off-axis source; that field can have a nonzero
transverse component, carried by the $m=1$ mode.

\section{Local image subtraction near the surface}

As a source charge approaches the conductor, the potential at the boundary becomes sharply
peaked.  If the charge is a distance $d$ from the surface, the angular width of
that peak scales with $d$, so the number of Fourier modes needed to resolve it
grows roughly like $1/d$, and brute-force mode
summation becomes expensive.

The solver therefore removes the most singular local behavior analytically.
Let $\mathbf{x}_0$ be the real charge position, and let $\mathbf{x}_0^\ast$ be
its reflection through the tangent plane at the nearest point on the conductor.
Place an auxiliary charge $-qw$ at $\mathbf{x}_0^\ast$, where $0\le w\le1$ is
a smooth cutoff that fades the correction away as the source moves farther from
the surface.  Call its potential $\Phi_{\mathrm{img}}$.

Instead of asking the BEM to reproduce the full sharply peaked boundary data,
solve only for the residual surface charge using the modified boundary data
$-(g_{\mathrm{direct}}+g_{\mathrm{img}})$.  The physical potential is then
reconstructed as
\begin{equation}
  \Phi_{\mathrm{induced}}
  = \mathcal{S}[\sigma_{\mathrm{res}}] + \Phi_{\mathrm{img}},
  \label{eq:reconstruct}
\end{equation}
where $\mathcal{S}$ denotes the surface-charge potential in
Eq.~\eqref{eq:single-layer}.

Provided the reflected charge lies inside the metal, the cutoff $w$ only
changes how the solution is split between the analytic image term and the BEM
residual.  It can change convergence dramatically, but it cannot change a
fully converged answer.

If the reflected point lands back in the vacuum, the logic above no longer
works, and the method fails.  This failure is most likely where the tangent plane is a bad local description
of the conductor, especially near a corner, edge, or cusp where the normal is
discontinuous or undefined, and can be avoided by smoothing the geometry.

\section{The infinite plane boundary}

We are interested in compact nano-features that are attached to an
otherwise infinite grounded plane.  Of course, the full infinite plane cannot be meshed
literally, and there are two ways to workaround this.  If the nano-feature protrudes above the plane, symmetry can be used to remove the plane exactly, and only a mesh on the feature itself is needed.  On the other hand, for a cavity cut into the plane, the same trick does not work and the infinite plane must instead be meshed out to some finite distance.

\subsection{Protrusions: remove the plane by mirror antisymmetry}

Assume the conductor fills the half-space below $z=0$, with a bounded body $B$
protruding into the upper half-space.  Reflect $B$ through $z=0$ and join it to
its mirror image.  Call the resulting closed, doubled body $\widetilde{B}$.

Now consider an excitation that is antisymmetric across the plane,
\begin{equation}
  \Phi_{\mathrm{ext}}(\rho,\phi,-z)
  = -\Phi_{\mathrm{ext}}(\rho,\phi,z).
  \label{eq:antisym}
\end{equation}
The doubled geometry is symmetric under $z\mapsto-z$, while the excitation
changes sign.  Therefore the induced solution must also change sign under that
reflection.  A function that is antisymmetric across $z=0$ is exactly zero on
$z=0$.

This means the grounded-plane boundary condition is satisfied automatically.
We can solve the exterior problem around the closed body $\widetilde{B}$ with no
plane in the mesh at all, then keep only the solution in the upper half-space.
There is no approximation in this step: the infinite plane has been removed by
symmetry, not truncated.

Both excitations used by the solver can be put in this antisymmetric form.

\begin{itemize}
\item \textbf{Uniform applied field.} The potential
  $\Phi_{\mathrm{ext}}=-E_z z$ is already antisymmetric in $z$.

\item \textbf{Point charge.} A single charge above the plane is not
  antisymmetric by itself.  Add an image charge of opposite sign at the
  reflected position below the plane.  The real-plus-image pair is
  antisymmetric and can therefore be used with the doubled body.
\end{itemize}

\subsection{Depressions: why the mirror trick no longer works}

Now reverse the geometry.  The conductor still fills the half-space below
$z=0$, but a cavity $D$ is cut into it.  The vacuum consists of the entire upper
half-space plus the interior of that cavity.

For a protrusion, reflecting the vacuum through $z=0$ gives a complementary
region below the plane.  The original and reflected vacuum regions fit together
without overlapping, and together they form exactly the space outside the
doubled closed body.  That is why the mirror construction can replace the plane
by an ordinary closed-surface BEM problem.

For a cavity, part of the vacuum already extends below the plane.  Reflecting
that cavity produces a region above the plane, but the upper half-space was
already vacuum.  The original and reflected vacuum regions therefore overlap
through a finite volume instead of fitting together as complementary halves.
There is no closed doubled body whose exterior is obtained by joining the two
copies.

For a depression, the plane must be included explicitly and cut off at a finite
radius.  That finite-radius cutoff is a genuine numerical approximation, so its
error must be checked by increasing the truncation radius until the result
converges.  This is a real qualitative difference between a bump and a cavity:
one admits an exact symmetry reduction, while the other requires a controlled
large-domain approximation.

\section{Superposition: several charges and independent groups}

The per-mode equations in Eq.~\eqref{eq:per-mode} are linear.  For several
charges, their Fourier boundary data can therefore be added first and solved in
one combined BEM solve.  The resulting induced surface charge is exactly the
sum of the responses that would have been obtained by solving each source
separately.

This is useful because every real charge then feels the field produced by the
same jointly polarized conductor.  The conductor-mediated interaction between
charges is therefore obtained naturally from a shared solve, without an extra
solve for every pair.

\section{Summary of the solver structure}

\begin{itemize}
\item The BEM first replaces the three-dimensional vacuum problem by a surface
  charge problem.  For a surface of revolution, a Fourier expansion in azimuth
  then reduces that surface problem to independent one-dimensional integral
  equations along the generating curve.

\item The per-mode kernel is the ring kernel in Eq.~\eqref{eq:ring-kernel},
  expressed through the toroidal harmonic $Q_{m-1/2}$.  There is no azimuthal
  mesh; source azimuths are represented analytically, and the only angular
  approximation is truncation at a finite $m=M$.

\item Truncation removes angular detail but preserves rotational covariance
  exactly because rotations mix only the cosine and sine components within a
  fixed $m$.  Absolute orientation therefore does not matter numerically; only
  relative azimuths do.

\item A nearby source requires many modes because its induced surface charge is
  sharply localized.  Subtracting a local tangent-plane image removes that
  singular part analytically and accelerates convergence.  The subtraction is
  exact as long as the reflected auxiliary charge lies inside the conductor.

\item A protrusion on an infinite grounded plane can be replaced exactly by a
  mirror-doubled body with an antisymmetric excitation.  For point charges, the
  required image source also supplies the physical image field of the plane.

\item A depression or cavity cannot use the same mirror reduction because the
  reflected vacuum overlaps the original vacuum rather than complementing it.
  The plane must then be truncated explicitly, and convergence with truncation
  radius must be checked numerically.
\end{itemize}

\end{document}
